\subsection{Inferenza e selezione delle mosse}

Durante una partita, il modello viene utilizzato per individuare la mossa da eseguire a partire dallo stato corrente della scacchiera. Questa fase è gestita dalla classe \texttt{TRMAgent}, implementata nel file \texttt{agents.py}.

A differenza dell'addestramento, durante l'inferenza i parametri della rete rimangono invariati. L'agente carica il modello, prepara l'input, esegue il forward pass e trasforma le predizioni in un'azione compatibile con le regole del gioco.

\subsubsection{Caricamento del modello e preparazione dell'input}

All'inizializzazione, l'agente carica il checkpoint \texttt{best\_model.pt}, dal quale recupera la configurazione necessaria a ricostruire il modello e i relativi pesi:

\begin{lstlisting}[language=Python, caption={Caricamento del modello dal checkpoint}, label={lst:inference_load}]
checkpoint = torch.load(checkpoint_path, map_location=self.device)
config_dict = checkpoint["config_dict"]

self.model = GoTRMModel(config_dict).to(self.device)
self.model.load_state_dict(checkpoint["model_state_dict"])
self.model.eval()
\end{lstlisting}

Il modello viene trasferito sul dispositivo disponibile, selezionato tra CUDA, Apple MPS e CPU, e impostato in modalità valutazione.

Prima dell'inferenza, il metodo \texttt{\_extract\_features()} costruisce il tensore di input a partire dallo stato mantenuto dall'ambiente. La rappresentazione è coerente con quella utilizzata durante il training e presenta forma \([1,7,N,N]\).

L'esecuzione del modello avviene senza calcolo dei gradienti:

\begin{lstlisting}[language=Python, caption={Esecuzione del forward pass durante l'inferenza}, label={lst:inference_forward}]
x = self._extract_features(env)

with torch.inference_mode():
    logits = self.model(x).squeeze(0)
\end{lstlisting}

Il risultato è un vettore di \(N^2+1\) logits, corrispondenti alle intersezioni e all'azione \textit{PASS}.

\subsubsection{Filtraggio delle mosse illegali}

Prima della selezione, è necessario escludere le azioni non consentite dalle regole del Go. L'agente richiede all'ambiente una maschera mediante il metodo \texttt{get\_legal\_moves\_mask()}, che restituisce un vettore booleano della stessa dimensione dei logits.

Le azioni illegali vengono escluse sostituendo i rispettivi valori con \(-\infty\):

\begin{lstlisting}[language=Python, caption={Mascheramento delle mosse illegali}, label={lst:legal_mask}]
masked_logits = logits.masked_fill(
    ~legal_mask,
    float("-inf")
)
\end{lstlisting}

In questo modo, dopo la normalizzazione, tali azioni ricevono probabilità nulla e non possono essere selezionate. La verifica rimane delegata all'ambiente, che costituisce l'unico riferimento per la legalità delle mosse.

\subsubsection{Penalizzazione della mossa PASS}

Per ridurre la tendenza del modello a interrompere prematuramente la partita, durante l'inferenza può essere applicata una penalizzazione al logit associato a \textit{PASS}.

La configurazione predefinita prevede una penalità pari a \(3.5\), applicata quando il numero di pietre presenti sulla scacchiera è inferiore a 16:

\begin{lstlisting}[language=Python, caption={Penalizzazione del logit PASS}, label={lst:pass_penalty}]
if occupied_stones < before_stones:
    adjusted_logits = logits.clone()
    adjusted_logits[-1] -= penalty
\end{lstlisting}

La penalizzazione interviene esclusivamente sul valore di \textit{PASS}, non modifica i pesi del modello e può essere regolata o disabilitata attraverso i parametri di configurazione.

\subsubsection{Conversione in probabilità e selezione della mossa}

Dopo il filtraggio e l'eventuale penalizzazione, i logits vengono convertiti in probabilità mediante softmax:

\[
    P(a_i) =
    \frac{\exp(l_i)}
    {\sum_{j \in A_{\mathrm{legali}}}\exp(l_j)}
\]

dove \(A_{\mathrm{legali}}\) rappresenta l'insieme delle azioni consentite.

La selezione non consiste necessariamente nell'individuare l'azione con probabilità massima. L'agente considera infatti tutte le mosse la cui probabilità risulta sufficientemente vicina a quella della migliore candidata.

Indicando con \(p_{\max}\) la probabilità massima, viene definita una soglia pari al 98\% di tale valore:

\[
    p_{\mathrm{soglia}}=0.98\cdot p_{\max}
\]

L'insieme delle mosse candidate è quindi:

\[
    A_{\mathrm{candidate}} =
    \left\{
    a_i \in A_{\mathrm{legali}}
    \mid P(a_i)\geq 0.98\cdot p_{\max}
    \right\}
\]

La mossa viene infine estratta uniformemente da questo insieme:

\[
    a_{\mathrm{scelta}}\sim
    \operatorname{Uniforme}(A_{\mathrm{candidate}})
\]

La tolleranza del 2\% è relativa alla probabilità massima e non rappresenta una differenza assoluta di due punti percentuali. In presenza di più azioni con probabilità simili, l'agente introduce quindi una componente stocastica, evitando di selezionare sistematicamente la stessa mossa.

La procedura è implementata nella funzione \texttt{select\_near\_best\_action()}. L'azione restituita viene infine passata all'ambiente, che ne gestisce l'esecuzione e l'aggiornamento dello stato della partita.
